\documentclass[11pt,a4paper]{article}

\usepackage[T1]{fontenc}
\usepackage[utf8]{inputenc}
\usepackage[american]{babel}
\usepackage[a4paper,margin=2.8cm]{geometry}
\usepackage{microtype}
\usepackage[hidelinks]{hyperref}

\hypersetup{
  pdftitle={That Summer at Fuller Lodge},
  pdfauthor={J.J. Gómez-Cadenas}
}

\setlength{\parindent}{1.25em}
\setlength{\parskip}{0.35em}
\raggedbottom

\title{That Summer at Fuller Lodge}
\author{J.J. Gómez-Cadenas}
\date{}

\begin{document}

\maketitle


Probably Emil Konopinski would not have remembered the flying-saucer cartoon from \emph{The New Yorker}, had it not been Monday. Any other day of the week, Kate was the most efficient of the Fuller Lodge girls, and the prettiest, with her short blonde hair and her smiling, bright red lips.

But it was Monday, and on Mondays, Kate was a walking ruin. Pale, sluggish, her makeup looking cheap, Kate moved around between the tables like a zombie, unable to cope with the orders. The heat---that summer of 1950 had broken every record of the decade---and the patrons’ hurried mood didn’t help either. It had been more than fifteen minutes since Emil and his group had taken their usual place at the restaurant, and Kate had not yet brought them the menus.

Fermi was starting to show signs of impatience, and Emil knew that the risk of him losing his temper and leaving before the girl managed to serve them lunch was increasing every second. And if Fermi left, Kate would be in serious trouble. She could even lose her job. Perhaps that would teach her not to show up for work in such a sorry state. But, on the other hand, except on Mondays, Kate was so nice, particularly to him. Emil really did appreciate the fact that she never forgot to refill his cup over and over. “Have some more coffee, Doctor Konoski,” she would say, always getting his name wrong, always smiling. Teller had also noticed that the girl was especially kind to him, and he always winked and said, “Look how the young lady cares for you, Doctor Casanova.” But Teller’s patience was as short as Fermi’s, and crossing him was even more dangerous.

One of Emil’s favorite private jokes---being classified material, it had to be private---was to liken Fermi to a neutrino, as fast as light in his mental calculations, and indifferent to everything but work. He only interacted with reality every now and then, but on those rare occasions he could explode like the atomic bomb they had built together. But Teller was, even on his good days, like the hydrogen weapon he was determined to build. If it could be built, it would be a thousand times more dangerous than the deadly bombs they had dropped on Hiroshima and Nagasaki.

Emil preferred not to think too much about that. He had seen the pictures, of course, and had been briefed, like everybody else, on the reasons. He knew that all that destruction had, probably, saved many lives in the end… And yet, there was this voice inside his head, repeating stubbornly, wrong, wrong, wrong.

And so, Emil decided to do something to help Kate and mute the voice. He reached into his pocket and produced the \emph{New Yorker} cartoon that he had cut out carefully from the magazine that morning.

“Look,” he said, handing the picture to the others. “Pretty cool, don’t you think?”

The cartoon showed a bunch of Martians, each wearing a little hat topped with a small antenna. The little green men---looking more like gnomes in pajamas than true aliens---had, apparently, landed in the city, and were busy carrying all the trash cans they could find to their flying saucer. The picture illustrated a sarcastic article in which the columnist accused the aliens of stealing the trash cans that had been disappearing from the streets of New York all summer.

“Nuts!” said Teller, bursting into one of his explosive laughs, while he punched Emil in the shoulder. “Journalists are just nuts!”

Fermi examined the drawing with the serious concentration that was his trademark. Emil could almost hear the wheels, spinning so fast, in his wonderful brain.

“A war!” roared Teller, his Hungarian accent as thick as ever. “We have just survived a fucking war and those damned journalists have nothing better to do than talk about Martians?”

“A \emph{New Yorker} joke,” offered Emil. Meanwhile, he noticed that Kate, alarmed by Teller’s bellowing, was waking up and hurrying toward their table, menus in hand.

“Sure,” agreed Teller. “But half of the city’s papers had taken the bait and are wasting ink speculating whether flying saucers are already here, quoting every crackpot who has seen little green men. Flying saucers! A Russian nuke can hit us at any moment and people have nothing better to do than discuss flying saucers.”

“It is a reasonable theory,” said Fermi.

Teller slammed his huge fist down on the table, sending a small quake through the ice in the large glasses of water Kate had just brought.

“What do you mean, Enrico? Is it reasonable that the newspapers worry about the Martians, rather than the Russians? Is it reasonable that the summer’s hot topic is the fate of New York’s trash cans?”

“The model is pretty good,” insisted Fermi, unperturbed, pointing to the cartoon with a tanned finger. “It explains both the disappearing trash cans and the saucer sightings.”

“Seriously, guys,” broke in Herbert York, who had been quiet so far, too busy reading a paper---and scratching on it with his red pen---to pay them any attention. “What if they were really here?”

“Nonsense!” said Teller. “The distance between stars is huge. Light takes four years to reach us from Proxima Centauri, our closest neighbor. At what speed can one of your flying saucers travel, Herbert? Ten percent of the speed of light seems too much already. So, forty years to get here from the next corner of our neighborhood. What is the distance to the galaxy’s center, twenty thousand light years, perhaps? Little green men would need two hundred thousand years to come from there. Seems like quite a long journey just to steal a few trash cans, don’t you think?”

“This guy claims that special relativity is wrong and light speed is not a limit,” said York, mercilessly slapping the paper he had been marking up.

“Another?” said Teller. “How many this year?”

“Twenty,” said York, shrugging.

“And how many were wrong?”

“All of them,” said York happily. “Including this one.”

“Then?” said Teller, opening his arms theatrically.

“Well, who knows,” said York. “Perhaps aliens live much longer than us and can afford space travel. Or perhaps they hibernate.”

“Robots,” said Teller. “They would probably send robots to steal our trash cans. Maybe they need the metal. What do you think, Enrico?”

“Think?” Fermi was distracted, chasing one of his bright ideas, uncoupled, as usual, from the rest of the world. “Think about what?”

“A little green man stealing trash cans,” said Teller.

“Oh. About a ten percent probability, I would say.”

Everyone laughed. Only Fermi would give the impossible a ten percent chance.

“There you go.” Teller blew his nose into the napkin that Kate had just put in front of him, together with the menu. “I’ll have a burger, honey.”

“Same for me,” said Emil, feeling much better. The crisis was over, and Kate would survive one more Monday.

“And for me,” said Herbert.

Fermi ordered black coffee and a tiny sandwich. He was eating less and less, Emil noticed. He had the chilling intuition that perhaps it would not be long before the great man decoupled for good and became a true neutrino.

“We’ve got enough trouble with the Russians,” said Teller, without letting go of his obsessions. “No need for Martians too, thank you.”

“Maybe they are friendly,” said Herbert.

“No way,” answered Teller, with a sardonic smile on his strangely feminine lips. “Anyone with space-travel technology would also have much better weapons than ours. Better if they can’t come. Because if they do, we are all toast.”

“Maybe not,” Emil dared to say, now that Kate was safely off to the counter with their orders. “Perhaps a civilization more advanced than ours would also be morally superior.”

“Bullshit,” cried Teller, punching him again, not very lightly, in the shoulder. “Do you think that we are morally superior to the Russians? Do you think they don’t love their children, or their dogs? They are good guys, the Russkies, most of them anyway. So tell me. Why in hell are we building our little toy? Why in hell are we trying to build the Super if they are so nice?”

“Not so loud, Edward,” York cautioned.

Teller looked around him, almost as if he expected to be surrounded by Soviet spies. The Super was the name of the H-bomb they hoped would give them an edge over the Russians. Or so the generals, and Edward Teller, believed. If it worked, it could deliver a thousand Hiroshimas in a single artifact, enough to erase Moscow from the face of the Earth. For the second time that day, Emil felt a chill running down his spine.

“Anyway, then,” said Teller, in a lower voice. “Russians or Martians, same story. Morality is not important here. Megatons are all that count, and any alien who can get here is bound to have more nukes than us. I’m just happy that light travels so slowly, guys.”

Kate arrived with the orders. In spite of the deep lines running under her eyes, and the lousy makeup, she smiled, smiled as always, while she poured coffee into his cup. Probably, Emil thought, somewhere in Russia, a pretty girl like her was smiling at a physicist who also loved his children, and also worked on a hydrogen bomb capable of erasing New York from the face of the Earth.

“Like her, don’t you, Casanova,” said Teller, with his mouth full of meat.

Before Emil had a chance to reply, Fermi set his cup of coffee down on the table---he had not touched his sandwich yet---and said:

“But where is everybody?”

Amazing, thought Emil, the way that even Teller would shut up when Fermi fired off one of his torpedoes. Everyone at the table knew that he was not joking. Once again, his mind had reached further than everyone else’s.

“Let’s see.” Fermi took his pen and started shooting numbers across the paper napkin. “Let’s say that there are about a hundred billion stars in the galaxy. Certo?”

It was inevitable, when Fermi was thinking deeply, that he leaked a few words in Italian. The pen spilled numbers, which spread, along with the ink, all over the porous tissue of the napkin.

“How many stars will have a habitable planet?”

“Say ten percent,” answered Teller, who had dropped his burger and was as absorbed in the argument as Fermi. The man could be nasty, but his talent was a close match to that of the master.

“Ten percent, giusto, giusto. So we have 10 billion planets good for life. In how many of these does life actually appear?”

“Perhaps in all of them,” said Emil. “Sooner or later, some chemist will find a way to test Oparin’s theories in the lab.”

“I agree,” said Teller. “Oparin was right, in spite of being Russian. Give me a good planet, with plenty of water, hydrogen, methane and ammonia, and life is unavoidable.”

“I don’t see it this way,” said York. “On the contrary, I am convinced that life is a rare miracle. Perhaps unique to this planet.”

“A miracle,” said Fermi, with one of his shy smiles. “Ten percent then. We’ve got a billion planets with life. In how many does intelligence appear?”

“In none,” said Teller, shaking his large head, pessimistically. “Including this one.”

“If life is a miracle, intelligence is an even greater one,” agreed York.

“Ten percent then, va bene. We are left with a hundred million planets where intelligence blossoms. Of these, maybe ten percent develop an advanced civilization. Of those, another ten percent would have developed long before us, say one to a hundred million years earlier. Vero? Thus, a million very advanced civilizations, which could have invented space travel a million years ago. Even if they travel at ten percent of light speed, as Edward suggests, they have had plenty of time to explore and colonize the whole galaxy. Certo, ragazzi?”

“I see your point, Enrico,” said Teller, without a trace of irony in his voice. “They should already be here.”

“Exactly. We should have seen them long ago. Or at least we should have heard them.”

“Radio signals,” said Emil. “Of course. The galaxy should be full of radio traffic.”

“Perhaps it’s true that they’re stealing the trash cans after all,” offered York, with a smile. He was the only one in the group unmoved by the argument. His beliefs were too strong to let him take the matter seriously.

“Not a joke, Herbert,” said Teller, dryly. “Enrico just nailed it. Where is everybody? The figures are very clear. Even if you give or take a couple of orders of magnitude, the result is still the same. We should have already seen the fucking Martians, unless they are hiding. And why should they hide from us, if they are technologically superior? I don’t get it.”

“I wouldn’t lose too much sleep thinking about this,” said York.

“Perhaps you should, Herbert. Ignorance can kill you.”

“Sometimes ignorance may be bliss,” argued York, with a priestly smile.

“Really? Ask around Hiroshima and Nagasaki about ignorance. Perhaps they were also convinced that nothing bad could happen to them, right up until the moment Little Boy and Fat Man hit them.”

“Can’t we leave Hiroshima and Nagasaki alone?” said York, looking hurt. “We were talking about…”

“We were talking about a serious question,” Teller cut in. “We were trying to understand why a galaxy that should be buzzing with signals from other civilizations is silent. We were trying to understand why aliens seem to be hiding. I can come up with a number of answers. None of them is good for us. And you know what? If everybody is silent, perhaps we should be quiet too. But we are just making more and more noise, sending radio waves to the galaxy, telling everyone where we are. Not smart, dear boy. Not smart.”

Third chill, thought Emil, and this one ran all the way to his gut, freezing the burger he had just eaten. Emil took a deep breath, trying to control the sudden nausea.

As if feeling his anxiety, Kate approached, brandishing the coffee pot.

“More coffee, Doctor Konoski?”

Emil sighed, grateful, and held out his cup toward her smile, painted in red.

\end{document}
